\documentclass{article}
\usepackage{geometry}
%\usepackage[showframe]{geometry}       % shows frame
%\usepackage[showframe,pass]{geometry}  % shows frame but does not change page geometry

%\usepackage{ametsoc}
\usepackage[intlimits]{amsmath}
\usepackage{bm}
%\usepackage{amssymb,latexsym,graphics}
\usepackage{natbib}  % all options handled by \bibstyle@XXX below

\usepackage{graphicx}
\usepackage{caption}
\usepackage{subfig}
%\usepackage{endfloat} % uncomment to push figures and tables to the end of the document

\usepackage{booktabs}
\usepackage[notcite,notref]{showkeys}
\usepackage[osf]{mathpazo} % palatino fonts % uncomment for Palatino font
\usepackage{todonotes}
%\usepackage[disable]{todonotes} % remove [disable] to see to-do notes
\definecolor{darkgreen}{rgb}{0,0.5,0}
\usepackage[linkcolor=blue,citecolor=darkgreen,colorlinks=true]{hyperref}

\usepackage{siunitx}

% + custom commands - - - - - - - - - - - - - - - - - - - - - - - - - - - - - - - - - -
\newcommand{\pderiv}[2]{\frac{\partial #1}{\partial #2}} % partial derivative
\newcommand{\vonkarman}{\kappa} % von Karman constant
\newcommand{\pt}{\theta} % potential temperature
\newcommand{\ustar}{\ensuremath{u_*}} % velocity scale
\newcommand{\bstar}{\ensuremath{b_*}} % buoyancy scale
\newcommand{\tstar}{\ensuremath{\pt_*}} % buoyancy scale
\newcommand{\rcrit}{\text{Ri}_c} % potential temperature

% typical panel width in multi-panel figures
\newlength{\panelW}
\setlength{\panelW}{7cm}

\captionsetup[subfigure]{position=top,textfont=scriptsize,labelfont=scriptsize,singlelinecheck=false,captionskip=0pt,farskip=2pt,margin=15pt}

\title{Exploration of Monin-Obukhov Similarity Theory and its implementation in GFDL models}

% + main body - - - - - - - - - - - - - - - - - - - - - - - - - - - - - - - - - - - - -
\begin{document}
\maketitle

% - - - - - - - - - - - - - - - - - - - - - - - - - - - - - - - - - - - - - - - - - - - -
\section{Basic equations of motion}

Following \cite{Monin1971}, start with the basic equations of motion in Boussinesq approximation:
\begin{subequations}
\begin{align} \label{eq:boussinesq}
   \pderiv{\bm{u}}{t} + (\bm{u} \cdot \nabla) \bm{u} & = -\frac{1}{\rho_0}\nabla{p'} + \nu \nabla^2 \bm{u} - \frac{\bm{g}}{p_0}\pt' \\
   \pderiv{\pt}{t} + (\bm{u} \cdot \nabla) \bm{\pt} & =  \chi\nabla^2 \pt
\end{align}
\end{subequations}
and continuity equation in approximation of incompressible liquid
\begin{equation}
   \nabla \cdot \bm{u} = 0
\end{equation}
where $\pt$ is potential temperature, $\pt'$ is the departure from mean, and $p'$ is the departure from hydrostatics.
Viscosity terms in the right-hand side are typically neglected.

Fluxes of momentum and heat within constant flux layer
\begin{gather}
   \tau = \rho_0 (\overline{u'w'}) = \text{const}_1 \\
   H = \rho_0 c_p (\overline{\pt'w'}) = \text{const}_2
\end{gather}

Friction velocity scale $\ustar$ is defined by the momentum sink at the surface,
\begin{equation}
   \tau = \rho_0 (\overline{u'w'}) = \rho_0 \ustar^2
\end{equation}

Similarly, temperature scale $\tstar$ is defined by the relationship
\begin{equation}
   H = \rho_0 c_p (\overline{\pt'w'}) = \rho_0 c_p \tstar \ustar
\end{equation}

% - - - - - - - - - - - - - - - - - - - - - - - - - - - - - - - - - - - - - - - - - - - -
\section{Dimensional reasoning}

\subsection{Neutral case}
In the neutral case, far from the wall (where the viscosity effects can be important) the vertical gradient of the mean wind can only depend on stress $\tau$, air density $\rho$, and vertical coordinate $z$ \cite[see][section 5.3]{Monin1971}, therefore the equation becomes:
\begin{equation}
   \pderiv{u}{z} = \frac{\ustar}{\vonkarman z}
\end{equation}

\subsection{Stratified flow}
The statistics of the flow in the near-surface layer can only depend on small number of parameters that appear in the above section, namely on temperature-to-buoyancy scale $g/\pt_0$, density $\rho_0$, fluxes of momentum $\tau$ (or, equivalently, velocity scale $\ustar$) and temperature $H/(\rho_0 c_p)=\overline{\pt'w'}$, and $z$. Units of these quantities are
\begin{equation}
    \left[\frac{g}{\pt_0}\right] = \si{\per\K\m\per\s\squared}; \quad
    \left[\rho_0\right] = \si{\kg\per\m\cubed}; \quad
    \left[\ustar\right] = \si{\m\per\s}; \quad
    \left[\overline{\pt'w'}\right] = \si{\K\m\per\s}
\end{equation}

There is only one way (to the arbitrary power) to combine the above parameters to construct the values of length dimension:
\begin{equation}
    \left[\frac{g}{\pt_0}\right]^{n_1}
    \left[\rho_0\right]^{n_2}
    \left[\ustar\right]^{n_3}
    \left[\overline{\pt'w'}\right]^{n_4} = [z]
\end{equation}
and that is
\begin{equation}
   L = \frac{\ustar^3}{\vonkarman\frac{g}{\pt_0}(\overline{\pt'w'})}
\end{equation}

Per \cite{Monin1971}, the dimensionless von Karman constant $\vonkarman$ was first introduced into the equation for $L$ by Obukhov in 1946 and has usually been preserved in all subsequent works by tradition. $L>0$ for stable situations; $L<0$ for unstable cases.

In non-neutral case,
\begin{equation}\label{eq:u-profile}
   \frac{\vonkarman z}{\ustar}\pderiv{u}{z} = \phi_m(\zeta)
\end{equation}
where
$\zeta=z/L$,
$\phi_m$ is stability correction function,
$\phi_m \to 1$ as $\zeta \to 0$ ro retrieve neutral case.

\cite{Held2004} then defines
\begin{equation}
   F_m(\zeta) = \int\frac{\phi_m}{\zeta}\,d \zeta
\end{equation}
to solve \eqref{eq:u-profile} as
\begin{equation}
   u(z) = \frac{\ustar}{\vonkarman} \left[
          F_m\left(\frac{z}{L}\right)-F_m\left(\frac{z_0}{L}\right)
      \right]
\end{equation}

Other literature, including \cite{Jimenez2012} use equivalent approach to define
\begin{equation}
   \psi_m(\zeta) = \int\frac{1-\phi_m(\zeta)}{\zeta}\, d \zeta
               = \ln(\zeta) - F_m(\zeta)
\end{equation}
to get
\begin{equation}
   u(z) = \frac{\ustar}{\vonkarman}\left[
          \ln\left(\frac{z}{z_0}\right) - \psi_m\left(\frac{z}{L}\right) + \psi_m\left(\frac{z_0}{L}\right)
      \right]
\end{equation}

Similarly, for buoyancy
\begin{equation}\label{eq:b-profile}
   \frac{\vonkarman z}{\bstar}\pderiv{b}{z} = \phi_b(\zeta)
\end{equation}
and
\begin{equation}
   b(z) - b(0) = \frac{\bstar}{\vonkarman} \left[
          F_b\left(\frac{z}{L}\right)-F_b\left(\frac{z_b}{L}\right)
      \right]
\end{equation}
or
\begin{equation}
   b(z) - b(0) = \frac{\bstar}{\vonkarman}\left[
          \ln\left(\frac{z}{z_b}\right) - \psi_h\left(\frac{z}{L}\right) + \psi_h\left(\frac{z_b}{L}\right)
      \right]
\end{equation}

\subsection{Turbulent diffusion coefficients}
By definition of turbulent diffusion coefficient
\begin{equation}
   (\overline{u'w'}) = K_m \pderiv{u}{z} = \ustar^2
\end{equation}
and substituting ${\partial u/\partial z}$ from~\eqref{eq:u-profile}
\begin{equation}
   K_m = \frac{\ustar^2}{\partial u/\partial z} = \frac{\ustar \vonkarman z}{\phi(\zeta)}
\end{equation}
% - - - - - - - - - - - - - - - - - - - - - - - - - - - - - - - - - - - - - - - - - - - -
\section{Stability functions}
\subsection{\protect\cite{Held2004}}
Stability functions recommended by Garratt.

\subsubsection{Unstable ($\zeta < 0$)}
\begin{align}
   \phi_m(\zeta) & = (1-16 \zeta)^{-1/4} \\
   \phi_b(\zeta) & = (1-16 \zeta)^{-1/2}
\end{align}

\subsubsection{Stable ($\zeta > 0$)}
\begin{equation}
    \phi_m(\zeta) = \phi_b(\zeta) = 1+5\zeta
\end{equation}
which are considered to have empirical support for $-5 < \zeta < 1$.

% - - - - - - - - - - - - - - - - - - - - - - - - - - - - - - - - - - - - - - - - - - - -
\section{Free convection limit}

As $u \to 0$, $\ustar \to 0$, and $\zeta \to -\infty$. In this case, vertical profile of temperature may become ill defined, depending on the asymptotic properties of the stability functions. Consequently, vertical flux of buoyancy $B$ becomes ill defined.

To see what is going on, start from equation for vertical profile of buoyancy \eqref{eq:b-profile}:
\begin{equation*}
    \frac{\vonkarman z}{\bstar}\pderiv{b}{z} = \phi_b(\zeta)
\end{equation*}
Also, from definition of \bstar~\eqref{eq:bstar}
\begin{equation*}
    \bstar = \frac{B}{\rho \ustar}
\end{equation*}
and definition of $\zeta$
\begin{equation*}
    \zeta = \frac{z}{L} = -\frac{z \vonkarman \bstar}{\ustar^2}
\end{equation*}
we get
\begin{equation}
    \pderiv{b}{z} = \frac{\bstar}{\vonkarman z}\phi_b(\zeta)
    = \frac{\bstar}{\vonkarman z}\phi_b\left(-\frac{z \vonkarman \bstar}{\ustar^2}\right)
    = \frac{B}{\vonkarman z \rho \ustar }\phi_b
      \left(-\frac{z \vonkarman B}{\rho \ustar^3}\right)
\end{equation}

When the wind is close to zero in non-neutral situation, it is reasonable to assume that the heat flux should not depend on $\ustar$. Therefore, in equation~\eqref{eq:b-profile}, the dependence on $\ustar$ should disappear when wind goes to zero (and, therefore, $\ustar \to 0$ and $\zeta \to \infty$). Suppose $\phi(\zeta) \propto \zeta^n$ for $\zeta \to -\infty$, with some power $n$. Then
\begin{equation}
    \frac{B}{\vonkarman z \rho \ustar }\phi_b \left(-\frac{z \vonkarman B}{\rho \ustar^3}\right)
  = \frac{B}{\vonkarman z \rho \ustar } C \left(-\frac{z \vonkarman B}{\rho \ustar^3}\right)^n
  = - C \frac{B^{n+1} z^{n-1} \vonkarman^{n-1}}{\rho^{n+1}} \frac{1}{\ustar^{3n+1}}
\end{equation}
For this expression to be finite and non-zero when $\ustar \to 0$, we should have $n=-1/3$.

For $n < -1/3$, the vertical gradient of buoyancy would tend to zero for finite heat flux $B$, which implies that the heat flux would tend to infinity for finite values of vertical gradient. On the other hand, if $n > -1/3$, the vertical gradient would tend to infinity for finite $B$.

\cite{Held2004} and \cite{Beljaars1994} propose to use ``gustiness'' to solve this problem: adding an ortogonal component to the atmospheric wind speed, determined by the $\bstar$, $\ustar$, and imposed parameters of the convective boundary layer:
\begin{gather}
   |U| = U^2 + V^2 + G^2 \\
   G = \beta \left(z_i \bstar \ustar \right)^{1/3}
\end{gather}
where $\beta$ is an empirical constant on the order of 1, and
$z_i$ is the characteristic height of the convective boundary layer \citep{Beljaars1994}. In GFDL models, $z_i$ is assumed to be constant \SI{1000}{m}.

This treatment should ensure the same asymptotic behavior as imposing the above constraint on $\phi_b$ at $\ustar \to -\infty$. However, such approach can only work if calculation of gustiness is embedded within the solution of the implicit equations of heat and momentum balances in the Monin-Obukhov layer. In GFDL model that is not done: rather, gustiness is calculated in the atmosphere based on the flux and stability values from the previous time step. Therefore, I believe there is no guarantee that the surface fluxes produce reasonable solution in convective limit.
% - - - - - - - - - - - - - - - - - - - - - - - - - - - - - - - - - - - - - - - - - - - -
\section{FMS}
Current FMS code (2020.04) has a number of options.

\subsection{Unstable ($\zeta < 0$)}
\begin{align}
   \phi_m(\zeta) & = (1-16 \zeta)^{-1/4} \\
   \phi_b(\zeta) & = \begin{cases}
      (1-16\zeta)^{-1/3}, & \text{for new option (Ming Zhao)} \\
      (1-16\zeta)^{-1/2}, & \text{for old option}
   \end{cases}
\end{align}
Note that the momentum stability function is the same as in \cite{Held2004}

\subsection{Stable ($\zeta > 0$), option 1}
\begin{equation}
   \phi_m(\zeta) = \phi_b(\zeta) = 1 + \left(5 + \frac{\zeta}{\rcrit}\right)\frac{\zeta}{1+\zeta} \\
\end{equation}

\subsection{Stable ($\zeta > 0$), option 2}
\begin{equation}
   \phi_m(\zeta) = \phi_b(\zeta) = \begin{cases}
      1+5\zeta, & \zeta < \zeta_t \\
      1+5 \zeta_t + \frac{\zeta-\zeta_t}{\rcrit}, & \zeta \ge \zeta_t
   \end{cases}
\end{equation}

Note that in stable case, $\phi_m(\zeta) = \phi_b(\zeta)$, like they were in \cite{Held2004}.

\subsubsection{Model settings}
ESM4.1 uses stable option 2 with $\rcrit=1.0$ and $\zeta_t=0.5$; it does not use new Ming's unstable option. Figures~\ref{fig:small-z0-sens} and \ref{fig:large-z0-sens} show typical fluxes resulting from Monin-Obukhov calculations in this case. Note that for the purposes of these calculations gustiness was not added to the atmospheric wind velocity.

One notable feature of these plots is how the sensible heat flux does not depend monotonically on the wind speed for the case of $k B^{-1}=0$ (i.e., when we assume $z_0=z_b$). This feature is even more striking in Figure~\ref{fig:sens-vs-u}, which shows results for high value of roughness length $z_0=\SI{2}{m}$, consistent with the values that land model reaches in forests.

Speaking of Figure~\ref{fig:sens-vs-u}, one cannot help noticing huge magnitudes of heat fluxes, reaching values of \SIrange{500}{1000}{\W\per\m\squared} even for modest temperature gradient of \SI{2}{\K} across the near-surface layer. Note how for slow wind speeds the heat flux rapidly shoots upwards. In both cases ($k B^{-1}=2$ and $k B^{-1}=2$) gustiness would have introduced significant changes in the heat fluxes, if it was applied.

\begin{figure}
    \centering
    \subfloat[ESM4.1, $k B^{-1}=2$]{
        \includegraphics[width=\panelW]{graphics/z0m0.1kb2/flux_t.pdf}
    }
    \subfloat[ESM4.1, $k B^{-1}=0$]{
        \includegraphics[width=\panelW]{graphics/z0m0.1kb0/flux_t.pdf}
    }

    \subfloat[Brutsaert, $k B^{-1}=2$]{
        \includegraphics[width=\panelW]{brutsaert/z0m0.1kb2/flux_t.pdf}
    }
    \subfloat[Brutsaert, $k B^{-1}=0$]{
        \includegraphics[width=\panelW]{brutsaert/z0m0.1kb0/flux_t.pdf}
    }
    \caption{Dependence of sensible heat flux on surface temperature, for roughness $z_0=\SI{0.1}{\m}$. Air temperature at the top of surface layer is kept at \SI{300}{\K}.}
    \label{fig:small-z0-sens}
\end{figure}

\begin{figure}
    \centering
    \subfloat[ESM4.1, $k B^{-1}=2$]{
        \includegraphics[width=\panelW]{graphics/z0m0.1kb2/b_star.pdf}
    }
    \subfloat[ESM4.1, $k B^{-1}=0$]{
        \includegraphics[width=\panelW]{graphics/z0m0.1kb0/b_star.pdf}
    }

    \subfloat[Brutsaert, $k B^{-1}=2$]{
        \includegraphics[width=\panelW]{brutsaert/z0m0.1kb2/b_star.pdf}
    }
    \subfloat[Brutsaert, $k B^{-1}=0$]{
        \includegraphics[width=\panelW]{brutsaert/z0m0.1kb0/b_star.pdf}
    }
    \caption{Dependence of $\bstar$ on surface temperature, for roughness $z_0=\SI{0.1}{\m}$.}
\end{figure}

\begin{figure}
    \centering
    \subfloat[ESM4.1, $k B^{-1}=2$]{
        \includegraphics[width=\panelW]{graphics/z0m0.1kb2/flux_m.pdf}
    }
    \subfloat[ESM4.1, $k B^{-1}=0$]{
        \includegraphics[width=\panelW]{graphics/z0m0.1kb0/flux_m.pdf}
    }

    \subfloat[Brutsaert, $k B^{-1}=2$]{
        \includegraphics[width=\panelW]{brutsaert/z0m0.1kb2/flux_m.pdf}
    }
    \subfloat[Brutsaert, $k B^{-1}=0$]{
        \includegraphics[width=\panelW]{brutsaert/z0m0.1kb0/flux_m.pdf}
    }
    \caption{Dependence of surface stress on surface temperature, for roughness $z_0=\SI{0.1}{\m}$.}
\end{figure}

\begin{figure}
    \centering
    \subfloat[ESM4.1, $k B^{-1}=2$]{
        \includegraphics[width=\panelW]{graphics/z0m0.1kb2/u_star.pdf}
    }
    \subfloat[ESM4.1, $k B^{-1}=0$]{
        \includegraphics[width=\panelW]{graphics/z0m0.1kb0/u_star.pdf}
    }

    \subfloat[Brutsaert, $k B^{-1}=2$]{
        \includegraphics[width=\panelW]{brutsaert/z0m0.1kb2/u_star.pdf}
    }
    \subfloat[Brutsaert, $k B^{-1}=0$]{
        \includegraphics[width=\panelW]{brutsaert/z0m0.1kb0/u_star.pdf}
    }
    \caption{Dependence of $\ustar$ on surface temperature, for roughness $z_0=\SI{0.1}{\m}$.}
\end{figure}

% - - - - - - - - - - - - - - - - - - - - - - - - - - - - - - - - - - - - - - - - - - - -
\begin{figure}
    \centering
    \subfloat[ESM4.1, $k B^{-1}=2$]{
        \includegraphics[width=\panelW]{graphics/z0m2kb2/flux_t.pdf}
    }
    \subfloat[ESM4.1, $k B^{-1}=0$]{
        \includegraphics[width=\panelW]{graphics/z0m2kb0/flux_t.pdf}
    }

    \subfloat[Brutsaert, $k B^{-1}=2$]{
        \includegraphics[width=\panelW]{brutsaert/z0m2kb2/flux_t.pdf}
    }
    \subfloat[Brutsaert, $k B^{-1}=0$]{
        \includegraphics[width=\panelW]{brutsaert/z0m2kb0/flux_t.pdf}
    }
    \caption{Dependence of sensible heat flux on surface temperature, for high value of roughness $z_0=\SI{2}{\m}$.}
    \label{fig:large-z0-sens}
\end{figure}

\begin{figure}
    \centering
    \subfloat[ESM4.1, $k B^{-1}=2$]{
        \includegraphics[width=\panelW]{graphics/z0m2kb2/b_star.pdf}
    }
    \subfloat[ESM4.1, $k B^{-1}=0$]{
        \includegraphics[width=\panelW]{graphics/z0m2kb0/b_star.pdf}
    }

    \subfloat[Brutsaert, $k B^{-1}=2$]{
        \includegraphics[width=\panelW]{brutsaert/z0m2kb2/b_star.pdf}
    }
    \subfloat[Brutsaert, $k B^{-1}=0$]{
        \includegraphics[width=\panelW]{brutsaert/z0m2kb0/b_star.pdf}
    }
    \caption{Dependence of $\bstar$ on surface temperature, for high value of roughness $z_0=\SI{2}{\m}$.}
\end{figure}

\begin{figure}
    \centering
    \subfloat[ESM4.1, $k B^{-1}=2$]{
        \includegraphics[width=\panelW]{graphics/z0m2kb2/flux_m.pdf}
    }
    \subfloat[ESM4.1, $k B^{-1}=0$]{
        \includegraphics[width=\panelW]{graphics/z0m2kb0/flux_m.pdf}
    }

    \subfloat[Brutsaert, $k B^{-1}=2$]{
        \includegraphics[width=\panelW]{brutsaert/z0m2kb2/flux_m.pdf}
    }
    \subfloat[Brutsaert, $k B^{-1}=0$]{
        \includegraphics[width=\panelW]{brutsaert/z0m2kb0/flux_m.pdf}
    }
    \caption{Dependence of surface stress on surface temperature, for high value of roughness $z_0=\SI{2}{\m}$.}
\end{figure}

\begin{figure}
    \centering
    \subfloat[ESM4.1, $k B^{-1}=2$]{
        \includegraphics[width=\panelW]{graphics/z0m2kb2/u_star.pdf}
    }
    \subfloat[ESM4.1, $k B^{-1}=0$]{
        \includegraphics[width=\panelW]{graphics/z0m2kb0/u_star.pdf}
    }

    \subfloat[Brutsaert, $k B^{-1}=2$]{
        \includegraphics[width=\panelW]{brutsaert/z0m2kb2/u_star.pdf}
    }
    \subfloat[Brutsaert, $k B^{-1}=0$]{
        \includegraphics[width=\panelW]{brutsaert/z0m2kb0/u_star.pdf}
    }
    \caption{Dependence of $\ustar$ on surface temperature, for high value of roughness $z_0=\SI{2}{\m}$.}
\end{figure}

% - - - - - - - - - - - - - - - - - - - - - - - - - - - - - - - - - - - - - - - - - - - -

\begin{figure}
    \centering
    \subfloat[ESM4.1, $k B^{-1}=2$]{
        \includegraphics[width=\panelW]{graphics/z0m2kb2-vsU/flux_t.pdf}
    }
    \subfloat[ESM4.1, $k B^{-1}=0$]{
        \includegraphics[width=\panelW]{graphics/z0m2kb0-vsU/flux_t.pdf}
    }

    \subfloat[Brutsaert, $k B^{-1}=2$]{
        \includegraphics[width=\panelW]{brutsaert/z0m2kb2-vsU/flux_t.pdf}
    }
    \subfloat[Brutsaert, $k B^{-1}=0$]{
        \includegraphics[width=\panelW]{brutsaert/z0m2kb0-vsU/flux_t.pdf}
    }

    \caption{Dependence of sensible heat flux on wind speed, for high value of roughness $z_0=\SI{2}{\m}$.}
    \label{fig:sens-vs-u}
\end{figure}

\begin{figure}
    \centering
    \subfloat[ESM4.1, $k B^{-1}=2$]{
        \includegraphics[width=\panelW]{graphics/z0m2kb2-vsU/b_star.pdf}
    }
    \subfloat[ESM4.1, $k B^{-1}=0$]{
        \includegraphics[width=\panelW]{graphics/z0m2kb0-vsU/b_star.pdf}
    }

    \subfloat[Brutsaert, $k B^{-1}=2$]{
        \includegraphics[width=\panelW]{brutsaert/z0m2kb2-vsU/b_star.pdf}
    }
    \subfloat[Brutsaert, $k B^{-1}=0$]{
        \includegraphics[width=\panelW]{brutsaert/z0m2kb0-vsU/b_star.pdf}
    }

    \caption{Dependence of $\bstar$ on wind speed, for high value of roughness $z_0=\SI{2}{\m}$.}
\end{figure}

\begin{figure}
    \centering
    \subfloat[ESM4.1, $k B^{-1}=2$]{
        \includegraphics[width=\panelW]{graphics/z0m2kb2-vsU/flux_m.pdf}
    }
    \subfloat[ESM4.1, $k B^{-1}=0$]{
        \includegraphics[width=\panelW]{graphics/z0m2kb0-vsU/flux_m.pdf}
    }

    \subfloat[Brutsaert, $k B^{-1}=2$]{
        \includegraphics[width=\panelW]{brutsaert/z0m2kb2-vsU/flux_m.pdf}
    }
    \subfloat[Brutsaert, $k B^{-1}=0$]{
        \includegraphics[width=\panelW]{brutsaert/z0m2kb0-vsU/flux_m.pdf}
    }

    \caption{Dependence of surface stress on wind speed, for high value of roughness $z_0=\SI{2}{\m}$.}
\end{figure}

\begin{figure}
    \centering
    \subfloat[ESM4.1, $k B^{-1}=2$]{
        \includegraphics[width=\panelW]{graphics/z0m2kb2-vsU/u_star.pdf}
    }
    \subfloat[ESM4.1, $k B^{-1}=0$]{
        \includegraphics[width=\panelW]{graphics/z0m2kb0-vsU/u_star.pdf}
    }

    \subfloat[Brutsaert, $k B^{-1}=2$]{
        \includegraphics[width=\panelW]{brutsaert/z0m2kb2-vsU/u_star.pdf}
    }
    \subfloat[Brutsaert, $k B^{-1}=0$]{
        \includegraphics[width=\panelW]{brutsaert/z0m2kb0-vsU/u_star.pdf}
    }

    \caption{Dependence of $\ustar$ on wind speed, for high value of roughness $z_0=\SI{2}{\m}$.}
\end{figure}

% - - - - - - - - - - - - - - - - - - - - - - - - - - - - - - - - - - - - - - - - - - - -
\bibliographystyle{ametsoc}
\bibliography{/Users/malyshev/work/PapersDB}


\end{document}